\section{Open-weight judges}
\label{sec:openweight}

The frontier study has three limits: closed models whose probabilities we
cannot read, vendors that may watermark their text, and few genres. We
repeat it with four smaller open-weight models, which avoid all three but are
much weaker judges, so this study says little about frontier models directly.

\subsection{Setup}

\paragraph{Panel and texts.} Qwen2.5-7B (Alibaba), Llama-3.1-8B (Meta),
Mistral-7B-v0.3 (Mistral AI) and Phi-4-mini (Microsoft), each run in 4-bit on
a single laptop GPU (Appendix~\ref{app:openweight}). With four candidates,
chance is 25\%. Each model answered the same 120 prompts as the API texts
(temperature 0.7), giving 480 texts, plus a second, independent answer to
every prompt. We generated the texts ourselves from public weights, so none
carries a watermark.

\paragraph{Formats.} The judges use the lineup and one-text formats with
their own four names, and answer ``Did you write this text?'' about every
text in a fresh conversation. Here we read the probability of ``yes'' from
the model, so we can compute an AUC. Two further conditions follow
\citet{khullar2026selfattribution}: the judge first answers the prompt
itself and is then asked the yes/no question about a candidate answer
(Section~\ref{sec:onpolicy}). We also read each judge's likelihood of every
text, and had each judge rate every text's quality (Section~\ref{sec:mechanism}).

\subsection{Results}

\begin{table*}[t]
\centering\small
\setlength\tabcolsep{3.5pt}
\begin{tabular}{l rrrrr rrr rr}
\toprule
 & \multicolumn{5}{c}{Lineup} & \multicolumn{3}{c}{Yes/no, fresh context} & \multicolumn{2}{c}{Own likelihood} \\
\cmidrule(lr){2-6}\cmidrule(lr){7-9}\cmidrule(lr){10-11}
Judge & Self & Peer & FA & Self-adv. [95\% CI] & Hit$-$FA & Hit & FA & AUC & Pick-self & AUC \\
\midrule
Qwen & 56.7\% & 15.8\% & 49.2\% & +40.8$^{*}$ [+31.1, +50.6] & +7.5 & 51.7\% & 58.6\% & 0.45 & 100.0\% & 1.00 \\
Llama & 32.5\% & 22.2\% & 36.1\% & +10.3$^{*}$ [+1.9, +18.6] & $-$3.6 & 85.8\% & 85.6\% & 0.53 & 100.0\% & 1.00 \\
Mistral & 25.0\% & 20.8\% & 20.3\% & +4.2 [$-$6.1, +14.4] & +4.7 & 99.2\% & 99.4\% & 0.41$^{*}$ & 100.0\% & 1.00 \\
Phi & 20.8\% & 22.8\% & 22.5\% & $-$1.9 [$-$10.6, +6.9] & $-$1.7 & 4.2\% & 1.4\% & 0.53 & 100.0\% & 1.00 \\
\bottomrule
\end{tabular}
\caption{Open-weight judges, all four genres pooled (120 prompts). Lineup: self-naming, peer-baseline and false-alarm (FA) rates, the self-advantage with a prompt-bootstrap interval, and the hit minus false-alarm rate, which is positive only if a judge names itself more on its own text than on others'. Yes/no: ``Did you write this text?'' in a fresh context; hit and false-alarm rates at $P(\text{yes})>0.5$ and the AUC of $P(\text{yes})$ for own versus others' text (0.5 is chance). Own likelihood: the share of prompts on which the text the judge finds most likely, relative to the other judges, is its own (chance 25\%), and the AUC of that relative likelihood. $^{*}$Holm-corrected within-prompt permutation $p<0.05$ across the four judges.}
\label{tab:openweight}
\end{table*}


\paragraph{No open-weight judge recognises itself in the lineup.} Qwen and
Llama beat their peers (Table~\ref{tab:openweight}), but neither
discriminates. Qwen names itself on \OWQwenLineupSelf{} of its own texts, but
also on \OWQwenLineupFA{} of the texts it did not write; Llama names itself
\emph{more} often on others' texts than on its own. Both put their own name on
much of what they read, which is enough to beat peers who rarely name them.
Overall lineup accuracy is \OWLineupAcc{}, at chance, even though a style
classifier attributes these texts at \OWClfPooled{}.

\paragraph{One text at a time, and asked directly, each judge answers by
habit.} Shown one text at a time, Qwen names itself on \OWQwenSingleNamesSelf{}
of all texts and Llama on every text, whoever wrote them; all four judges have
a significant self-advantage and none discriminates
(Appendix~\ref{app:more}). Asked ``Did you write this text?'', Mistral says
yes to \OWMistralBinaryHit{} of its own texts and \OWMistralBinaryFA{} of
everyone else's, while Phi says yes to \OWPhiBinaryHit{} and
\OWPhiBinaryFA{}. Read as self-naming rates, these would rank the judges from
near-certain self-recognition to none. The AUC, which ignores how readily a
judge says yes, is at or below chance for all four (0.41 to 0.53); Mistral's
is reliably \emph{below} chance.

\subsection{Answering the task first}
\label{sec:onpolicy}

\begin{table}[t]
\centering\small
\setlength\tabcolsep{4pt}
\begin{adjustbox}{max width=\linewidth}
\begin{tabular}{l rrl rrl rrl}
\toprule
 & \multicolumn{3}{c}{Asked in a fresh conversation} & \multicolumn{3}{c}{After answering: same answer} & \multicolumn{3}{c}{After answering: a new answer} \\
\cmidrule(lr){2-4}\cmidrule(lr){5-7}\cmidrule(lr){8-10}
Judge & Hit & FA & AUC [95\% CI] & Hit & FA & AUC [95\% CI] & Hit & FA & AUC [95\% CI] \\
\midrule
Qwen & \OWQwenBinaryHit & \OWQwenBinaryFA & \OWQwenBinaryAUC{} \OWQwenBinaryAUCCI & \OWQwenOnPolHit & \OWQwenOnPolFA & \OWQwenOnPolAUC{} \OWQwenOnPolAUCCI & \OWQwenOnPolResHit & \OWQwenOnPolResFA & \OWQwenOnPolResAUC{} \OWQwenOnPolResAUCCI \\
Llama & \OWLlamaBinaryHit & \OWLlamaBinaryFA & \OWLlamaBinaryAUC{} \OWLlamaBinaryAUCCI & \OWLlamaOnPolHit & \OWLlamaOnPolFA & \OWLlamaOnPolAUC{} \OWLlamaOnPolAUCCI & \OWLlamaOnPolResHit & \OWLlamaOnPolResFA & \OWLlamaOnPolResAUC{} \OWLlamaOnPolResAUCCI \\
Mistral & \OWMistralBinaryHit & \OWMistralBinaryFA & \OWMistralBinaryAUC{} \OWMistralBinaryAUCCI & \OWMistralOnPolHit & \OWMistralOnPolFA & \OWMistralOnPolAUC{} \OWMistralOnPolAUCCI & \OWMistralOnPolResHit & \OWMistralOnPolResFA & \OWMistralOnPolResAUC{} \OWMistralOnPolResAUCCI \\
Phi & \OWPhiBinaryHit & \OWPhiBinaryFA & \OWPhiBinaryAUC{} \OWPhiBinaryAUCCI & \OWPhiOnPolHit & \OWPhiOnPolFA & \OWPhiOnPolAUC{} \OWPhiOnPolAUCCI & \OWPhiOnPolResHit & \OWPhiOnPolResFA & \OWPhiOnPolResAUC{} \OWPhiOnPolResAUCCI \\
\bottomrule
\end{tabular}
\end{adjustbox}
\caption{Open-weight judges asked ``Did you write this text?'' (120 own and
360 other texts per judge and condition). Left: in a fresh conversation.
Middle: after answering the prompt, about either that same answer or a peer's.
Right: after answering the prompt, about either a separate answer of its own
or a peer's. Hit and false-alarm rates count a probability of ``yes'' above
0.5; the AUC uses the probability itself (0.5 is chance).}
\label{tab:owyesno}
\end{table}


\citet{khullar2026selfattribution} find that a monitor's judgment shifts most
when it evaluates something it produced earlier in the same conversation.
Does that position help a judge recognise its writing? In the
\emph{verbatim} condition the judge answers the prompt and is then shown
either that same answer, word for word, or a peer's answer. Every judge's
AUC is now above chance (Table~\ref{tab:owyesno}), but this is a low bar,
since the candidate is a copy of text already in the conversation; Llama's
AUC is only \OWLlamaOnPolAUC{}. In the \emph{resampled} condition the judge's
own turn is a fresh answer and the candidate is its other answer or a peer's,
so recognising it requires recognising its style rather than a copy. Qwen
and Phi now tell their own answer from peers' (AUC \OWQwenOnPolResAUC{} and
\OWPhiOnPolResAUC{}), while Llama and Mistral fall reliably \emph{below}
chance (\OWLlamaOnPolResAUC{} and \OWMistralOnPolResAUC{}). Seeing its own
answer to the task helps two judges and misleads the other two.
